%! Simplifying Rational Expressions — Unit 1, Lesson 1.7 — Student Packet Cover Sheet
\documentclass[10pt]{article}
\usepackage{atda-article}
\usepackage{atda-boxes}
\usepackage{ltablex}
\keepXColumns

\begin{document}

% Full-bleed royal banner
\begin{tikzpicture}[remember picture, overlay]
  \fill[royal] (current page.north west)
    rectangle ([yshift=-0.9in]current page.north east);
\end{tikzpicture}
\vspace{-0.8in}

\begin{center}
  {\color{white}\textbf{\LARGE Algebra, Trigonometry, and Data Analysis}}\\[4pt]
  {\color{mistmid}\large\textbf{Unit 1: Algebraic Expressions and Factoring}}\\[6pt]
  {\color{white}\textbf{\large Lesson 1.7 \quad Simplifying Rational Expressions}}
\end{center}

\vspace{0.22in}
\namedateperiod
\vspace{0.18in}

\begin{learningtargetbox}
\small
\begin{enumerate}[leftmargin=1.5em, itemsep=5pt]
  \item \ldots{} find the \textbf{excluded values} of a rational expression by setting its
        denominator equal to zero and solving --- Lesson 1.6's method, with a new name.
  \item \ldots{} \textbf{simplify a rational expression} by factoring the top and the bottom and
        dividing out common \textbf{factors}, and explain why terms joined by $+$ or $-$ never
        cancel.
  \item \ldots{} \textbf{multiply, divide, add, and subtract} rational expressions, simplifying the
        result.
  \item \ldots{} \textbf{justify} that a simplified expression is equivalent to the original one ---
        for every value except the excluded ones --- and interpret a rational expression in context.
\end{enumerate}
\end{learningtargetbox}

\vspace{0.14in}

\begin{tocbox}
\small
\renewcommand{\arraystretch}{1.5}
\begin{tabularx}{\linewidth}{@{} c l X r @{}}
\rowcolor{mist}
\textbf{\#} & \textbf{Component} & \textbf{Description} & \textbf{Score} \\
\midrule
1 & Warm-Up        & Why $3+4$ over $3+5$ is not $4$ over $5$
                                                                    & \blank{1.2cm} \\
2 & Guided Notes   & Excluded values, simplifying, and the four operations
                                                                    & \blank{1.2cm} \\
3 & Group Activity & The backdrop: the area and the width are known, the length is not
                                                                    & \blank{1.2cm} \\
4 & Exit Ticket    & Two expressions and one answer to critique     & \blank{1.2cm} \\
5 & Homework       & Mixed practice, a garden bed, and a look ahead to Unit 2
                                                                    & \blank{1.2cm} \\
\midrule
  & \multicolumn{2}{r}{\textbf{Total}} & \blank{1.2cm} \\
\end{tabularx}
\end{tocbox}

\vspace{0.14in}

\begin{remindbox}
\small
A \textbf{rational expression} is just a fraction whose top and bottom are polynomials
(\texttt{A2.EO.1a--b}), and it obeys every rule ordinary fractions already obey. Two of those rules
run the whole lesson. First, \textbf{you may divide out only common \emph{factors}} --- things
being multiplied. Terms joined by $+$ or $-$ never cancel, which is why $\frac{x+3}{x+5}$ is
already as simple as it gets. Second, \textbf{you may never divide by zero}, so before you simplify
anything, find the values that make the denominator zero and carry them along as \textbf{excluded
values} --- read from the \emph{original} denominator, even after a factor cancels.
\end{remindbox}

\end{document}
